\documentclass[conference,compsoc]{IEEEtran}

\ifCLASSOPTIONcompsoc
  \usepackage[nocompress]{cite}
\else
  \usepackage{cite}
\fi

\usepackage{amsmath,amssymb}
\usepackage{graphicx}
\usepackage{booktabs}
\usepackage{braket}
\usepackage{xcolor}
\usepackage{multirow}
\usepackage{siunitx}

\hyphenation{multi-variate time-series cross-scheme}

\newcommand{\TODO}[1]{\textcolor{red}{\textbf{[TODO: #1]}}}

\begin{document}

\title{Quantum Probability Image Encoding for Multivariate\\
Time-Series Classification: Regime Diagnostics\\
Across Five UEA Benchmarks}

\author{%
\IEEEauthorblockN{Sogodam Atadana\IEEEauthorrefmark{1},
                  Özlem Karabiber Cura\IEEEauthorrefmark{2}}
\IEEEauthorblockA{%
\IEEEauthorrefmark{1}İzmir Katip Çelebi University,
Biomedical Engineering Department,
atadanas@gmail.com\\
\IEEEauthorrefmark{2}İzmir Katip Çelebi University,
Biomedical Engineering Department,
ozlem.karabiber@ikcu.edu.tr}}

\maketitle

% ==================================================================
\begin{abstract}
Quantum Probability Image Encoding (QPIE) was developed for spatial
multi-channel data such as RGB medical images. We test whether it
transfers without modification to multivariate time-series
classification (MTSC): each fixed-stride temporal patch is treated as
a pixel, the three accelerometer or pen-force channels replace the
three colour channels, and the five 3-qubit encoding schemes (Sep,
CRyE, GBE, CP-2L, CSE) and multi-basis Z$+$X$+$Y measurement pipeline
are reused unchanged.
Across five 3-channel UEA benchmarks --- the complete 3-dimensional
UEA subset --- the central result is negative.
On the four datasets where the encoding is functional, no entangling
scheme beats the separable baseline (Sep) at $p < 0.05$, and the
data-conditional scheme (CSE) falls significantly below Sep. The
image-domain channel-correlation mechanism that makes CSE useful on
RGB data does not operate on temporal sensor signals, whose
inter-channel correlations are sensor-physical rather than
class-conditional.
The pooled-to-CNN accuracy gap splits the datasets into regimes:
below 5 points where per-patch statistics carry the label
(Epilepsy, CharacterTrajectories), and 26--38 points where temporal
ordering does (UWaveGestureLibrary, Handwriting). On
EthanolConcentration both rungs collapse to near-random, a third
regime in which patch subsampling of the 1751-step series discards the
discriminative spectral detail.
On absolute accuracy QPIE is competitive in data-scarce settings but
reaches no new state of the art. It trails InceptionTime by 5.6 points
on the native UWave split and LITEMVTime by 2.9 points on
Epilepsy~\cite{ismail2025lite}; on Handwriting it exceeds every
deep-learning baseline in the LITE survey yet stays below the
strongest convolutional classifiers~\cite{ruiz2020bakeoff}.
Multivariate ROCKET run on the raw 3-channel series (same splits)
matches or beats QPIE on every benchmark except CharacterTrajectories.
With 896 training samples --- the same gesture instances stacked as
three channels (the \textit{UWaveGestureLibraryAll} benchmark) --- the
separable CNN ensemble reaches 97.5\%, within a point of multivariate
ROCKET on the raw 3-channel stack (98.4\%)~\cite{dempster2020rocket}.
\end{abstract}

\begin{IEEEkeywords}
quantum machine learning, time-series classification, QPIE,
multivariate UEA benchmarks, regime diagnostics, entanglement schemes.
\end{IEEEkeywords}

\IEEEpeerreviewmaketitle

% ==================================================================
\section{Introduction}
\label{sec:intro}

Multivariate time-series classification (MTSC) is central to
biomedical sensing: seizure detection from multi-lead EEG, gesture
recognition from wrist-worn inertial measurement units, and
handwriting analysis from pen-digitiser recordings all reduce to
labelling fixed-length windows of multi-channel signals.
The UEA archive~\cite{bagnall2018uea} provides standardised benchmarks
that a growing family of deep methods has attacked: InceptionTime,
ROCKET, ConvTran, and LITE set the current state of the art on most
problems~\cite{ismail2025lite,dempster2020rocket,foumani2024convtran}.

Quantum feature extractors remain largely unexplored in this space.
QPIE~\cite{yao2017quantum} maps an $n$-channel spatial vector into a
pure $n$-qubit state through angle encoding, then recovers features via
multi-basis projective measurement.
In a companion paper~\cite{atadana2026qpie} we applied QPIE to
3-channel RGB medical images (BloodMNIST, PathMNIST) and found that
(i) mean-pooling over pixels creates a permutation-invariance ceiling
that prevents fixed-unitary entanglement schemes from separating their
scores, and (ii) a data-conditional entanglement scheme (CSE) carries
a detectable image-domain signal that can be traced to inter-channel
pixel correlations ($\beta_\rho$ mechanism).

The present paper asks whether this framework, designed for spatial
data, transfers to temporal data without modification.
The structural analogy is direct: a time-series with $C$ channels
sampled over a window of $T$ timesteps is segmented into patches of
stride $s$, and each patch's $C$-dimensional sample is treated as one
pixel's channel vector.
The encoding circuits, multi-basis measurement, and aggregation
operators are reused exactly.

The transfer is not guaranteed to work.
Temporal datasets differ from images in three ways relevant to
quantum encoding:
(a) inter-channel correlations in IMU or pen-force data are driven by
    sensor physics, not by label-relevant structure;
(b) temporal ordering within a series is often class-discriminating
    in a way that mean-pooling over patches destroys; and
(c) training set sizes in UEA benchmarks are frequently below 200
    samples, where high-dimensional concat features overfit.

Our experiments expose all three effects, producing a diagnostic
framework: given a dataset, the ratio of pooled-features accuracy to
CNN accuracy predicts which effect dominates and therefore how much
of the quantum encoding's representational budget is actually used.

\noindent The contributions are:
\begin{enumerate}
  \item \textbf{Temporal adaptation of QPIE.}
    A direct, reusable mapping from MTSC patches to the QPIE
    encoding pipeline: same circuits, same observables, same
    aggregation operators.
  \item \textbf{No entanglement advantage on MTSC.}
    On the four datasets where the encoding is functional, no entangling
    scheme exceeds the separable baseline (Sep) at $p < 0.05$, and the
    data-conditional scheme (CSE) is significantly below Sep. This is
    consistent with the image-domain $\beta_\rho$ channel-correlation
    signal being specific to spatial colour structure rather than
    temporal sensor coupling. A control feeding the QPIE features to a
    strong classifier (ROCKET) confirms the encoding is lossy relative
    to the raw signal.
  \item \textbf{Regime diagnostic.}
    A dataset-level characterisation based on the pooled-to-CNN
    accuracy gap, observed across five benchmarks spanning three
    regimes (temporal-statistics, temporal-order, and encoding-failure).
  \item \textbf{Dual-split evaluation.}
    Benchmark results on canonical native splits for direct comparison
    to published baselines, alongside expanded-split results at
    deployment-scale training sizes.
\end{enumerate}

Section~\ref{sec:background} reviews QPIE and MTSC baselines.
Section~\ref{sec:method} describes the temporal adaptation.
Section~\ref{sec:experiments} states the experimental protocol.
Section~\ref{sec:results} reports results across all five benchmarks.
Section~\ref{sec:discussion} interprets the regime split and
discusses scope.

% ==================================================================
\section{Background}
\label{sec:background}

\subsection{QPIE and multi-basis measurement}

QPIE~\cite{yao2017quantum} assigns one qubit per spatial channel.
Given a normalised channel vector $\boldsymbol{\theta} \in [0, \pi/2]^n$,
angle encoding prepares
\begin{equation}
  \ket{\psi(\boldsymbol{\theta})} =
  \bigotimes_{c=1}^{n}
  \bigl(\cos\theta_c \ket{0} + \sin\theta_c \ket{1}\bigr)
  \label{eq:sep}
\end{equation}
for the separable baseline (Sep).
Entangling variants insert additional two-qubit gates before or after
the angle-encoding layer.
Multi-basis measurement projects each $n$-qubit state onto the
$\{0,1\}^n$ eigenbasis of the $Z$, $X$, and $Y$ Pauli operators in
turn, yielding $3 \cdot 2^n$ probability features per state.
For $n=3$ this gives 24 features per patch.
Our prior work~\cite{atadana2026qpie} fixed $n=3$ and compared six
encoding schemes across 14 image classes.

\subsection{Multivariate TSC baselines}

ROCKET~\cite{dempster2020rocket} applies a large bank of random
convolutional kernels and uses the resulting max-pooled features with
a ridge classifier.
InceptionTime~\cite{fawaz2020inceptiontime} stacks multiple inception
modules with multi-scale 1D convolutions.
LITE~\cite{ismail2025lite} is a lightweight variant optimised for
data-scarce and compute-constrained settings.
ConvTran~\cite{foumani2024convtran} combines 1D convolution with
transformer self-attention.
Ismail-Fawaz et al.~\cite{ismail2025lite} report a comprehensive
comparison across the full UEA archive; we draw all published baseline
numbers from that survey.

% ==================================================================
\section{Method}
\label{sec:method}

\subsection{From multivariate time series to quantum patches}
\label{sec:patching}

A multivariate series $\mathbf{X} \in \mathbb{R}^{C \times T}$
with $C = 3$ channels and $T$ timesteps is divided into patches of
stride $s$. Each patch position $p$ yields a channel vector
$\mathbf{x}_p \in \mathbb{R}^3$.
A per-channel robust scaler, fit on the training set, maps each
coordinate into $[0, 1]$ via the 5th-to-95th-percentile range.
The scaled vector is then converted to angles
$\theta_c = (\pi/2)\bar{x}_{p,c}$.

The analogy to QPIE image encoding is exact: a timestep with $C$
channels plays the role of a pixel with $C$ colour channels, and the
$P = \lfloor(T - 1)/s\rfloor + 1$ patches play the role of the
$H \times W$ pixels in an image.
No modification to the encoding circuits or measurement operators
is required.

\subsection{Encoding schemes}

We use the five encoding schemes that were validated on MedMNIST~\cite{atadana2026qpie}:
\begin{itemize}
  \item \textbf{Sep}: separable baseline (Eq.~\ref{eq:sep}).
  \item \textbf{CRyE}: controlled-$R_y$ entangler — each qubit is
    rotated by an angle conditioned on the previous qubit's state.
  \item \textbf{GBE}: global Bell entangler — a CNOT ladder
    applied after Sep to entangle all three qubits.
  \item \textbf{CP-2L}: two-layer circuit alternating angle
    encoding and CNOT entangling layers.
  \item \textbf{CSE}: data-conditional scheme — entangling rotation
    angles are shifted by scaled pairwise channel correlations
    $\rho_{ij}$ computed within each recording (across its patches),
    so the shift is a per-sample function of that recording's own
    inter-channel structure and uses no label or cross-sample
    information.
\end{itemize}

All schemes produce a 3-qubit state per patch.
Multi-basis Z$+$X$+$Y measurement yields 24 features per patch.

\subsection{Classification pipeline}

Three classification heads process the $P \times 24$ feature tensor.

\textbf{Pooled (L2-LR).}
Per-patch features are aggregated across patches with four marginal
moments --- mean, standard deviation, skewness, and kurtosis --- and
the 10th, 25th, 50th, 75th, and 90th percentiles, giving nine
univariate statistics per feature, together with the
upper-triangular cross-feature correlations.
With the $24$ per-patch features this is
$9 \times 24 = 216$ marginal statistics and
$\binom{24}{2} = 276$ correlation terms, a
$216 + 276 = 492$-dimensional recording-level descriptor.
A $\ell_2$-regularised logistic regression is trained on this fixed
vector.
This head is identical to the MedMNIST pipeline and isolates
pooling-extractable statistics.

\textbf{Concat-LR.}
Per-patch features are concatenated in temporal order into a single
vector of feature dimension $d = P \times 24$ (here from $360$ to
$5256$ across the five datasets at the coarse stride) and fed to the
same logistic regression. Here $d$ is this feature dimension and $N$
the number of training recordings (Table~\ref{tab:datasets}); because
$N$ is as small as $120$--$150$, the head operates in the $d \gg N$
regime and relies on $\ell_2$ regularisation to generalise.
It preserves temporal order at the cost of this high-dimensional
penalty.

\textbf{1D CNN.}
Per-patch features are treated as a sequence of $P$ frames, each of
dimension 24.
A three-block 1D convolutional network with 128 hidden units,
kernel size 3, and global average pooling followed by a linear
classifier is trained over 200 epochs with AdamW.
This head can exploit temporal ordering without the high-dimensional
penalty of Concat-LR.
A convolutional sequence head is preferred over a recurrent one
(LSTM/GRU): with $P \le 79$ frames and as few as $120$ training
recordings, the locality bias, translation invariance, and low
parameter count of a 1D CNN with global average pooling resist
overfitting and match the standard TSC baselines, whereas recurrent
models are data-hungry and underperform on UEA/UCR benchmarks at this
sample size.
An ensemble over 10 independently seeded training runs (soft-vote
over predicted class probabilities) is used as the headline result.

\textbf{Cross-scheme ensemble.}
Per-scheme CNN soft-vote distributions are averaged.
The 4-scheme ensemble excludes CSE;
the 5-scheme ensemble includes all five schemes.

% ==================================================================
\section{Experimental Setup}
\label{sec:experiments}

\subsection{Datasets}

Five 3-channel UEA datasets --- the complete set of 3-dimensional
problems in the UEA archive --- span the observed regimes and four
application domains (Table~\ref{tab:datasets}).

\begin{table}[t]
\caption{Dataset summary. Split sizes from the UEA archive.}
\label{tab:datasets}
\centering
\begin{tabular}{lrrrr}
\toprule
Dataset & $N_\text{tr}$ & $N_\text{te}$ & $K$ & $T$ \\
\midrule
Epilepsy             & 137  & 138  & 4  & 206 \\
CharacterTrajectories& 1422 & 1436 & 20 & 119 \\
UWaveGestureLibrary  & 120  & 320  & 8  & 315 \\
Handwriting          & 150  & 850  & 26 & 152 \\
EthanolConcentration & 261  & 263  & 4  & 1751 \\
\bottomrule
\end{tabular}
\end{table}

\textit{Epilepsy}~\cite{bagnall2018uea} contains 3-axis accelerometer
recordings of four activities (epileptic seizure simulation, walking,
running, and sawing), consistent with $K = 4$ in
Table~\ref{tab:datasets}.

\textit{CharacterTrajectories}~\cite{bagnall2018uea} captures
pen-digitiser data (x-velocity, y-velocity, pen-force) for 20
character classes.

\textit{UWaveGestureLibrary}~\cite{bagnall2018uea} provides
3-axis accelerometer data for 8 gesture classes, recorded from
the wrist.

\textit{Handwriting}~\cite{bagnall2018uea} contains 3-axis
pen-digitiser data for 26 uppercase letter classes.
With only 150 training samples across 26 classes, this is the
most data-scarce dataset in our benchmark.

\textit{EthanolConcentration}~\cite{bagnall2018uea} contains 3-channel
spectroscopic recordings ($T = 1751$) of distilled-spirit bottles at
four ethanol concentrations; the discriminative information lies in
fine spectral structure spread across the long series, and this
dataset functions as a stress test for the patch-subsampling encoding.

For UWaveGestureLibrary we run two evaluation protocols.
The \textit{native split} (120/320) is the UEA multivariate
benchmark used by Ismail-Fawaz et al.~\cite{ismail2025lite} and
allows direct SOTA comparison.
The \textit{stacked-univariate split} (896/3582) stacks the three
archival univariate axes (UWaveGestureLibraryX/Y/Z) as three channels,
reconstructing the 3-channel signal at 7.5$\times$ the training sample
count; labels are aligned and verified.
These are the same gesture instances as the UCR
\textit{UWaveGestureLibraryAll} benchmark, which concatenates the
identical three axes into a single univariate series. Published
bake-off baselines on UWaveGestureLibraryAll therefore provide a
direct, same-instance comparison, with the caveat that they operate on
the concatenated univariate form rather than our 3-channel stack.

\subsection{Covariate details and patch parameters}

Each patch is a single timestep sampled at stride $s$, so the patch
count is $P = \lfloor (T-1)/s \rfloor + 1$ and the stride sets temporal
resolution by subsampling rather than windowing.
The pooled and Concat-LR ablations use a coarse stride $s = 8$: the
pooled descriptor is fixed at $d = 492$ regardless of $P$ (it
aggregates across patches), whereas Concat-LR scales as $d = P \times
24$. The CNN uses a fine stride $s = 4$, yielding denser sequences of
up to $P = 79$ frames --- the maximum, on UWaveGestureLibrary with
$T = 315$ ($\lfloor 314/4 \rfloor + 1 = 79$) --- for the convolution
to exploit.
All per-channel scalers are fit on training data only and applied
to test data without leakage.
Ten independent seeds are used per configuration; paired $t$-tests
over the 10-seed distributions test scheme-vs-Sep comparisons.

\subsection{Baselines}

Published accuracy figures for InceptionTime, ROCKET,
ConvTran, D-CNN, FCN, ResNet, LITEMV, and LITEMVTime are taken from
Ismail-Fawaz et al.~\cite{ismail2025lite}, Table~6.
UWaveGestureLibraryAll baselines (ROCKET, HIVE-COTE, TS-CHIEF,
InceptionTime, ResNet) are the UCR ``bake off'' accuracies reported by
Dempster et al.~\cite{dempster2020rocket}, Table~1; the Handwriting
SOTA is the multivariate bake-off best of Ruiz et
al.~\cite{ruiz2020bakeoff}.
We additionally run multivariate ROCKET ourselves on the identical
native splits and the stacked UWave split (\texttt{aeon}
implementation of~\cite{dempster2020rocket}); these same-representation
numbers appear as ROCKET$^\S$ in Table~\ref{tab:main} and in
Table~\ref{tab:stacked}.

% ==================================================================
\section{Results}
\label{sec:results}

\subsection{Multi-dataset benchmark (native splits)}
\label{sec:main_results}

Table~\ref{tab:main} reports the headline CNN ensemble accuracy for
each dataset alongside the best published baseline.

\begin{table}[t]
\caption{CNN ensemble accuracy on native UEA splits.
Best QPIE scheme per dataset is named; deep-learning baselines from
Ismail-Fawaz et al.~\cite{ismail2025lite}.
$\Delta$ = QPIE best $-$ published SOTA (percentage points).
$^\dagger$Handwriting SOTA is the multivariate bake-off best
(InceptionTime~\cite{ruiz2020bakeoff}); the best LITE-survey
deep-learning result is 0.400 (LMVT), which QPIE exceeds by 13.9.
$^\S$Multivariate ROCKET~\cite{dempster2020rocket} on the raw
3-channel series and identical splits (this work).}
\label{tab:main}
\centering
\setlength{\tabcolsep}{4pt}
\begin{tabular}{lcccc}
\toprule
Dataset & Best QPIE & ROCKET$^\S$ & SOTA & $\Delta$ \\
\midrule
Epilepsy              & 0.964 (Sep)  & 0.986 & 0.993 (LMVT)  & $-2.9$  \\
CharacterTrajectories & 0.996 (GBE)  & 0.993 & 0.996 (LMVT)  & $+0.0$  \\
UWave (native)        & 0.850 (GBE)  & 0.941 & 0.909 (IT)    & $-5.6$  \\
Handwriting           & 0.539 (Sep)  & 0.587 & 0.657 (IT)$^\dagger$ & $-11.8$ \\
EthanolConcentration  & 0.274 (CP-2L)& 0.434 & 0.692 (LMVT)  & $-41.8$ \\
\bottomrule
\end{tabular}
\end{table}

On \textit{CharacterTrajectories} QPIE matches the LITEMVTime SOTA
(0.996 vs.\ 0.996), and on \textit{Epilepsy} it lands within 2.9
points of the 0.993 published best --- both with a single
109\,704-parameter CNN.
On \textit{UWaveGestureLibrary}, the 5-scheme cross-ensemble of 0.853
matches LITEMVTime (0.847) and ties LITEMV and FCN (0.850), falling
5.6 points below InceptionTime's 0.909.
On \textit{Handwriting}, the Sep CNN ensemble reaches 0.539. This
exceeds every deep-learning baseline in the LITE survey (best LMVT
0.400) but stays below the strongest convolutional classifiers in the
multivariate bake-off (InceptionTime 0.657, ResNet 0.598, ROCKET
0.567~\cite{ruiz2020bakeoff}). The 35-point spread between the two
baseline sets reflects how sensitive this 26-class, 150-sample problem
is to evaluation protocol; QPIE is competitive on it rather than
state of the art.

To remove any baseline-protocol ambiguity we ran multivariate
ROCKET~\cite{dempster2020rocket} on the raw 3-channel series and the
identical splits underlying QPIE's encoding (Table~\ref{tab:main},
ROCKET$^\S$ column). On equal footing ROCKET matches or beats the quantum encoding
on every dataset except CharacterTrajectories: Epilepsy $-2.2$, UWave
$-9.1$, Handwriting $-4.8$ points relative to QPIE. A fast classical
convolutional baseline is therefore at least as accurate as QPIE
throughout. The contribution of this work is consequently the transfer
behaviour, the regime diagnostic, and the entanglement-null, not
competitive accuracy.

A sharper test asks whether the encoding helps a strong classifier
rather than the logistic and CNN heads above. We fed each scheme's
QPIE per-patch features --- a 24-channel, length-$P$ series --- to the
same ROCKET, and compared against ROCKET on the raw 3-channel input
(Table~\ref{tab:rocketqpie}). The encoding is lossy: ROCKET on raw
equals or beats ROCKET on the best QPIE feature set on three of the
four functional datasets (Epilepsy $-2.2$, UWave $-1.9$, Handwriting
$-9.8$ points), tying only on the saturated CharacterTrajectories
($+0.2$). Entanglement again provides no benefit (GBE is the worst
input on UWave and Handwriting). The quantum feature map discards
label-relevant structure that a random-convolution baseline extracts
directly from the raw signal.

\begin{table}[t]
\caption{The QPIE encoding as a feature transform. ROCKET accuracy on
the raw 3-channel series vs.\ on the best-scheme QPIE per-patch
features (24-channel, length-$P$). Both use the same \texttt{aeon}
ROCKET; native splits.}
\label{tab:rocketqpie}
\centering
\setlength{\tabcolsep}{6pt}
\begin{tabular}{lcc}
\toprule
Dataset & ROCKET (raw) & ROCKET (best QPIE feat.) \\
\midrule
Epilepsy              & 0.986 & 0.964 (CRyE)  \\
CharacterTrajectories & 0.993 & 0.995 (CP-2L) \\
UWave (native)        & 0.941 & 0.922 (CSE)   \\
Handwriting           & 0.587 & 0.489 (CP-2L) \\
\bottomrule
\end{tabular}
\end{table}

The full ablation table (pooled, concat, CNN per scheme per dataset)
is given in Table~\ref{tab:ablation}.

\begin{table*}[t]
\caption{Accuracy across ablation rungs and encoding schemes.
Pooled: L2-LogReg on 492-dim aggregated features.
Concat-LR: L2-LogReg on temporally ordered patch features.
CNN ens: 10-seed soft-vote ensemble.
Values are mean over 10 seeds; CNN ens is the soft-vote accuracy.
Bold marks the numerically highest CNN-ens (soft-vote) accuracy per
dataset; per-seed CNN $\sigma$ ranges 0.003 (CharacterTrajectories) to
0.018 (UWave), so top-scheme differences at the CNN rung are within
seed noise (the only significant scheme effect is CSE\,$<$\,Sep,
Table~\ref{tab:cse}).}
\label{tab:ablation}
\centering
\setlength{\tabcolsep}{5pt}
\begin{tabular}{llccccc}
\toprule
Dataset & Rung & Sep & CRyE & GBE & CP-2L & CSE \\
\midrule
\multirow{3}{*}{Epilepsy}
  & Pooled    & $0.937 \pm 0.015$ & $0.912 \pm 0.019$ & $0.910 \pm 0.022$
               & $0.919 \pm 0.015$ & $0.899 \pm 0.015$ \\
  & Concat-LR & $0.820 \pm 0.028$ & $0.549 \pm 0.041$ & $0.622 \pm 0.019$
               & $0.744 \pm 0.021$ & $0.754 \pm 0.018$ \\
  & CNN ens   & $0.964$           & $0.942$           & $0.935$
               & $0.957$           & $0.942$           \\
\midrule
\multirow{3}{*}{CharacterTrajectories}
  & Pooled    & $0.960 \pm 0.004$ & $0.959 \pm 0.004$ & $0.956 \pm 0.004$
               & $0.956 \pm 0.004$ & $0.960 \pm 0.005$ \\
  & Concat-LR & $0.980 \pm 0.002$ & $0.981 \pm 0.002$ & $0.979 \pm 0.002$
               & $0.976 \pm 0.001$ & $0.981 \pm 0.002$ \\
  & CNN ens   & $0.995$           & $0.995$           & $\mathbf{0.996}$
               & $0.996$           & $0.995$           \\
\midrule
\multirow{3}{*}{UWave (native)}
  & Pooled    & $0.569 \pm 0.024$ & $0.525 \pm 0.030$ & $0.530 \pm 0.017$
               & $0.547 \pm 0.027$ & $0.553 \pm 0.031$ \\
  & Concat-LR & $0.834 \pm 0.022$ & $0.805 \pm 0.023$ & $0.679 \pm 0.016$
               & $0.812 \pm 0.019$ & $0.804 \pm 0.017$ \\
  & CNN ens   & $0.828$           & $0.828$           & $\mathbf{0.850}$
               & $0.828$           & $0.803$           \\
\midrule
\multirow{3}{*}{Handwriting}
  & Pooled    & $0.157 \pm 0.022$ & $0.134 \pm 0.012$ & $0.146 \pm 0.013$
               & $0.150 \pm 0.015$ & $0.118 \pm 0.014$ \\
  & Concat-LR & $0.218 \pm 0.029$ & $0.215 \pm 0.025$ & $0.179 \pm 0.013$
               & $0.210 \pm 0.022$ & $0.188 \pm 0.022$ \\
  & CNN ens   & $\mathbf{0.539}$  & $0.507$           & $0.404$
               & $0.492$           & $0.397$           \\
\midrule
\multirow{3}{*}{EthanolConcentration}
  & Pooled    & $0.283 \pm 0.024$ & $0.292 \pm 0.020$ & $0.269 \pm 0.026$
               & $0.265 \pm 0.014$ & $0.289 \pm 0.020$ \\
  & Concat-LR & $0.308 \pm 0.008$ & $0.321 \pm 0.017$ & $0.337 \pm 0.017$
               & $0.332 \pm 0.018$ & $0.309 \pm 0.023$ \\
  & CNN ens   & $0.255$           & $0.247$           & $0.236$
               & $\mathbf{0.274}$  & $0.221$           \\
\bottomrule
\end{tabular}
\end{table*}

\subsection{Regime diagnostic}
\label{sec:regime}

The pooled-to-CNN accuracy gap separates the five datasets into three
regimes (Table~\ref{tab:regime}).

\begin{table}[t]
\caption{Regime classification by pooled-to-CNN accuracy gap
(Sep scheme, native split).}
\label{tab:regime}
\centering
\begin{tabular}{lccc}
\toprule
Dataset & Pooled & CNN ens & Gap \\
\midrule
Epilepsy              & 0.937 & 0.964 & $+$2.7 \\
CharacterTrajectories & 0.960 & 0.995 & $+$3.5 \\
\midrule
UWave (native)        & 0.569 & 0.828 & $+$25.9 \\
Handwriting           & 0.157 & 0.539 & $+$38.2 \\
\midrule
EthanolConcentration  & 0.283 & 0.255 & $-$2.8 \\
\bottomrule
\end{tabular}
\end{table}

\textit{Epilepsy} and \textit{CharacterTrajectories} show small gaps
(2.7 and 3.5 points): class-discriminating information is concentrated
in per-patch amplitude and frequency statistics that mean-pooling
preserves.
These are the \textit{temporal-statistics} datasets.

\textit{UWaveGestureLibrary} and \textit{Handwriting} show large gaps
(25.9 and 38.2 points): temporal ordering is required, and pooling
destroys it.
These are the \textit{temporal-order} datasets.

\textit{EthanolConcentration} has a small gap ($-2.8$ points: the CNN
does not beat pooling) yet both rungs sit near the 25\% four-class
chance level. A small gap normally signals the temporal-statistics
regime; here it instead marks an \textit{encoding-failure} regime. The
discriminative signal is fine spectral detail spread along a 1751-step
series, and sampling one timestep every $s$ steps discards it before
encoding. The gap therefore predicts whether the CNN will beat
pooling, not whether the encoding captured usable signal at all ---
EthanolConcentration shows the two questions are distinct, and bounds
the diagnostic.
Here CP-2L is the one scheme to significantly edge Sep (by 1.7 points,
$p = 0.04$), but both sit near the 25\% four-class floor, so it is not
a meaningful advantage; the entanglement-null is therefore stated only
for the four functional datasets.

The Concat-LR rung does not resolve the gap on Handwriting or UWave
because $d \gg N_\text{train}$ (456-dimensional concatenation with
150 or 120 training samples) causes overfitting that $\ell_2$
regularisation only partially corrects.
The 1D CNN escapes this by treating patches as a short sequence
rather than a flat vector.

Mean inter-channel Pearson $\lvert\rho\rvert$ values
(Epilepsy 0.006, CharacterTrajectories 0.077, Handwriting 0.098,
UWave 0.201) do not predict the regime cluster:
UWave has the highest inter-channel correlation yet falls in the
temporal-order regime.
This confirms that the relevant structure is temporal order of patches,
not the spatial correlation between channels.

\subsection{CSE ablation: domain-specificity of the $\beta_\rho$ mechanism}
\label{sec:cse}

CSE is significantly below Sep on all five datasets at the CNN rung
(Table~\ref{tab:cse}).
The gaps are largest on temporal-order datasets (Handwriting $-$13.0
points, UWave $-$2.9 points), small but still significant on
temporal-statistics datasets (Epilepsy $-$0.9 points, $p = 0.013$;
CharacterTrajectories $-$0.1 points, $p = 0.009$), and $-$2.4 points
on the near-chance EthanolConcentration ($p < 0.001$).

\begin{table}[t]
\caption{CSE vs Sep at the CNN rung. $\Delta$ = CSE $-$ Sep mean
accuracy; $p$-value from paired $t$-test over 10 seeds.
$^*$ $p < 0.05$.}
\label{tab:cse}
\centering
\begin{tabular}{lccc}
\toprule
Dataset & $\Delta$ (ppts) & $p$ & Sig. \\
\midrule
Epilepsy              & $-0.87$ & 0.013 & $^*$ \\
CharacterTrajectories & $-0.08$ & 0.009 & $^*$ \\
UWave (native)        & $-2.91$ & 0.001 & $^*$ \\
Handwriting           & $-12.97$& $<$0.001 & $^*$ \\
EthanolConcentration  & $-2.43$ & $<$0.001 & $^*$ \\
\bottomrule
\end{tabular}
\end{table}

On PathMNIST (RGB images), CSE outperformed Sep by exploiting
image-domain colour correlations: the inter-pixel $\beta_\rho$ signal
is label-relevant because colour distributions are class-conditional.
On UEA time series, inter-channel correlations reflect sensor
coupling (e.g., gravity components in a wrist IMU) rather than class
membership.
The CSE scheme shifts entangling angles by these correlations and
thereby injects class-irrelevant structure into the quantum state,
degrading classification.
The deficit is stable across seeds: on Handwriting, where it is
largest, the per-seed standard deviation is 0.012, an order of
magnitude below the 13-point mean gap, so the drop reflects the
encoding rather than training variance. The fixed Bell entangler (GBE)
collapses by a similar margin on Handwriting (CNN mean 0.364),
indicating that on fine-grained order-dependent tasks global
entanglement degrades the per-patch features generally; CSE's
data-conditional shifts add nothing beyond this.

The cross-scheme ensemble also underperforms the best single-scheme
CNN ensemble on four of five datasets (Epilepsy $-0.7$, Handwriting
$-2.1$, CharacterTrajectories $-0.1$, EthanolConcentration $-3.8$
points): CSE degrades the ensemble when included.
Only on UWave (8 classes, some scheme diversity) does the cross-scheme
ensemble provide a small gain ($+0.3$ points over GBE alone).

\subsection{Practical performance: stacked-univariate UWave split}
\label{sec:utility}

On the stacked split (896 train, 3582 test) the Sep 1D CNN ensemble
reaches 97.5\%; the 5-scheme cross-ensemble adds no meaningful gain
(97.6\%, Table~\ref{tab:stacked}).
Multivariate ROCKET on the raw 3-channel stack reaches 98.4\%
(this work), so on equal footing QPIE trails ROCKET by 0.8 points; it
matches the univariate-concat bake-off ROCKET (97.5\%) and edges
HIVE-COTE and TS-CHIEF (96.9\%).

\begin{table}[t]
\caption{Accuracy on the \textit{UWaveGestureLibraryAll} benchmark
(896/3582) --- the same gesture instances as our stacked-univariate
split, in concatenated univariate form. Univariate baselines are the
UCR ``bake off'' accuracies from Dempster et al.~\cite{dempster2020rocket},
Table~1; ``ROCKET (3-channel, this work)'' is multivariate ROCKET on
our identical 3-channel stack.}
\label{tab:stacked}
\centering
\begin{tabular}{lc}
\toprule
Method & Accuracy \\
\midrule
ResNet                                   & 0.860 \\
InceptionTime~\cite{fawaz2020inceptiontime} & 0.955 \\
HIVE-COTE~\cite{bagnall2020hivecote}     & 0.969 \\
TS-CHIEF                                 & 0.969 \\
ROCKET~\cite{dempster2020rocket}         & 0.975 \\
\midrule
ROCKET (3-channel, this work)            & \textbf{0.984} \\
QPIE Sep (pooled, L2-LR)                 & 0.711 \\
QPIE Sep (concat, L2-LR)                 & 0.932 \\
QPIE Sep (CNN ens, 10 seeds)             & 0.975 \\
QPIE 5-scheme cross-ens                  & 0.976 \\
\bottomrule
\end{tabular}
\end{table}

The gap between the native split (85.3\%) and the stacked split
(97.6\%) reflects the training set size: 120 vs 896 samples.
The 7.5$\times$ increase in training data lifts QPIE from 5.6 points
below the native-split SOTA to parity with the strongest published
univariate baseline, though multivariate ROCKET on the same 3-channel
stack remains 0.8 points ahead.
The Table~\ref{tab:stacked} comparison uses UWaveGestureLibraryAll
bake-off numbers, which evaluate the identical instances in
concatenated univariate form; the only methodological difference is
that QPIE consumes the three axes as channels rather than as one
concatenated series.

% ==================================================================
\section{Discussion}
\label{sec:discussion}

\subsection{The pooled-to-CNN gap as a practical diagnostic}

The regime split is visible without running the full CNN ablation:
the pooled rung alone predicts whether the CNN will provide a large
lift.
Datasets where pooled accuracy is within 10 points of
a reasonable ceiling for the problem (e.g., Epilepsy at 93.7\% with
4 classes) are likely in the temporal-statistics regime, and the
pooled features can serve as a low-cost representation.
Datasets where pooled accuracy is far short of a usable level
(e.g., Handwriting at 15.7\% --- well above the 3.8\% random baseline
for 26 classes, but far below the CNN's 53.9\%) are in the
temporal-order regime, and the CNN is necessary.
The gap predicts whether the CNN will help, not whether the encoding
succeeds at all: on EthanolConcentration pooled and CNN agree near the
four-class chance level because the encoding never captured the
discriminative spectral detail.

This has a practical reading outside the quantum ML context: it
identifies whether a dataset's label information is concentrated in
local amplitude/frequency statistics (amenable to any mean-pooling
feature extractor) or in temporal ordering (requiring sequence models).
The quantum encoding provides a particular projection of the
per-patch distribution, but the regime split itself is a property of
the dataset and classifier family, not of the quantum circuit.

\subsection{Efficiency of pooled features}

On CharacterTrajectories (temporal-statistics regime) the Sep pooled
rung reaches 96.0\% with a 492-dimensional vector and an L2
logistic regression.
Wall-clock extraction time is under 2 seconds on a consumer laptop
(Apple M4, MPS backend).
The CNN adds 128 seconds per scheme but only 3.5 accuracy points
on this dataset.
For deployment in resource-constrained settings the pooled rung alone
is competitive.

\subsection{Limitations}

The 3-qubit architecture is fixed.
All circuits, observables, and aggregation operators are
simulation-based; no quantum hardware is used.
At $n = 3$ qubits the encoding is classically simulable: the schemes
are evaluated as fixed feature maps and the study makes no claim of
quantum computational advantage. The contribution is the transfer
behaviour and the regime characterisation, not a speedup.
The five benchmarks are the complete set of 3-dimensional problems in
the UEA archive; the $n$-qubit generalisation needed for
higher-channel datasets remains future work.
The regime diagnostic is empirical: the pooled-to-CNN gap threshold
is derived from five datasets and may not generalise to domains where
both effects are present simultaneously. EthanolConcentration already
exposes one failure mode --- a small gap that reflects encoding
failure rather than the temporal-statistics regime.

% ==================================================================
\section{Conclusion}
\label{sec:conclusion}

The QPIE 3-qubit encoding pipeline transfers from RGB medical
images to 3-channel multivariate time-series with no circuit
modifications.
Across five UEA benchmarks --- the complete 3-dimensional subset ---
the pooled-to-CNN accuracy gap separates temporal-statistics datasets
(small gap: Epilepsy, CharacterTrajectories) from temporal-order
datasets (large gap: UWave, Handwriting), with EthanolConcentration a
third case where both rungs fail near chance because patch subsampling
discards its spectral detail.
On the native UWave benchmark the cross-scheme ensemble (85.3\%)
ties lightweight baselines but trails the InceptionTime SOTA by
5.6 points.
On Handwriting the Sep CNN ensemble (53.9\%) tops every deep-learning
baseline in the LITE survey but trails the strongest bake-off
classifiers.
On the UWaveGestureLibraryAll split (896 training samples) the
separable CNN ensemble reaches 97.5\%, within a point of multivariate
ROCKET on the raw 3-channel stack (98.4\%).
The data-conditional entanglement scheme (CSE) is consistently below
Sep on all four functional datasets, consistent with the image-domain
$\beta_\rho$ mechanism being inactive when inter-channel correlations
are sensor-physical rather than label-conditional.
These results do not position QPIE as a competitive classifier: a
same-representation ROCKET matches or beats it on every dataset except
the saturated CharacterTrajectories, and feeding QPIE features into
ROCKET is lossy relative to the raw signal. The contribution is the
clean transfer of the encoding, the entanglement-null, and the regime
diagnostic --- a dataset-level prior for whether per-patch statistics,
temporal order, or neither carries the label.

% ==================================================================
\ifCLASSOPTIONcompsoc
  \section*{Acknowledgments}
\else
  \section*{Acknowledgment}
\fi

The authors declare that AI writing assistance was used in the
preparation of this manuscript. All quantitative claims were verified
against notebook outputs by the authors.

% ==================================================================
\begin{thebibliography}{99}

\bibitem{yao2017quantum}
X.-W.~Yao, H.~Wang, Z.~Liao, M.-C.~Chen \emph{et al.},
``Quantum image processing and its application to edge detection:
Theory and experiment,''
\emph{Phys. Rev. X}, vol.~7, p.~031041, 2017.

\bibitem{atadana2026qpie}
S.~Atadana and Ö.~Karabiber Cura,
``When does entanglement help in quantum image encoding?
A per-image mechanism test on {MedMNIST},''
in \emph{Proc. IEEE Int. Conf. Quantum Comput. Eng. (QCE)}, 2026,
to appear.

\bibitem{bagnall2018uea}
A.~Bagnall, H.~A.~Dau, J.~Lines, M.~Flynn, J.~Large, A.~Bostrom,
P.~Southam, and E.~Keogh,
``The {UEA} multivariate time series classification archive, 2018,''
\emph{arXiv:1811.00075}, 2018.

\bibitem{ismail2025lite}
A.~Ismail-Fawaz, M.~Devanne, S.~Berretti, J.~Weber, and G.~Forestier,
``Look into the {LITE} in deep learning for time series classification,''
\emph{Int. J. Data Sci. Anal.}, vol.~20, pp.~4029--4049, 2025.

\bibitem{dempster2020rocket}
A.~Dempster, F.~Petitjean, and G.~I.~Webb,
``{ROCKET}: Exceptionally fast and accurate time series classification
using random convolutional kernels,''
\emph{Data Min. Knowl. Discov.}, vol.~34, no.~5, pp.~1454--1495, 2020.

\bibitem{fawaz2020inceptiontime}
H.~Ismail-Fawaz, B.~Lucas, G.~Forestier, C.~Pelletier, D.~F.~Schmidt,
J.~Weber, G.~I.~Webb, L.~Idoumghar, P.-A.~Muller, and F.~Petitjean,
``{InceptionTime}: Finding {AlexNet} for time series classification,''
\emph{Data Min. Knowl. Discov.}, vol.~34, no.~6, pp.~1936--1962, 2020.

\bibitem{foumani2024convtran}
N.~M.~Foumani, C.~W.~Tan, G.~I.~Webb, and M.~Salehi,
``Improving position encoding of transformers for multivariate time
series classification,''
\emph{Data Min. Knowl. Discov.}, vol.~38, no.~1, pp.~22--48, 2024.

\bibitem{ruiz2020bakeoff}
A.~P.~Ruiz, M.~Flynn, J.~Large, M.~Middlehurst, and A.~Bagnall,
``The great multivariate time series classification bake off: a review
and experimental evaluation of recent algorithmic advances,''
\emph{Data Min. Knowl. Discov.}, vol.~35, no.~2, pp.~401--449, 2021.

\bibitem{bagnall2020hivecote}
A.~Bagnall, M.~Flynn, J.~Large, J.~Lines, and M.~Middlehurst,
``On the usage and performance of the hierarchical vote collective of
transformation-based ensembles version 1.0 ({HIVE-COTE} v1.0),''
in \emph{Advanced Analytics and Learning on Temporal Data (AALTD)},
ser. LNAI, vol.~12588.\ \ Springer, 2020, pp.~3--18.

\end{thebibliography}

\end{document}
